\responsetoreviewer{对审稿人 1 的回复}
\reviewersection{审稿人 1，意见 1（R1.1）}
\begin{reviewercomment}
分割与分类结果大多以图像层面报告。作者应澄清患者层面与检查层面的划分、跨数据集的去重，以及同一患者的多个帧或视图是否会跨划分出现。患者层面的再分析和聚类感知的置信区间会增强对报告性能的信心。
\end{reviewercomment}
\begin{authorresponse}
\begin{itemize}
\item 感谢审稿人。我们已按下列方式修订统计报告。
\item 独立性与近重复：在 \textbf{TN3K、TN5K、ThyroidXL、DDTI、PKTN、ZJH-8K} 等\textbf{各来源队列之间}做感知哈希（dHash，Hamming ≤ 5）筛查（见文末 \textbf{表 R1}）。其中与 \textbf{审稿人 2 --- 意见 4}（ZJH 与 TN3K 关系）最相关的是：TN3K（官方 train+val+test）与 ZJH-8K 共 \textbf{1156} 对 H=0 图像对；若仅看 TN3K \textbf{官方 test} 与 ZJH-8K，则为 \textbf{2} 对 H=0。
\item 聚类定义：ZJH-8K、ThyroidXL 使用可复现 patient ID；无 ID 队列 \textbf{一图一簇}。Bootstrap \textbf{1000} 次、\textbf{seed=42}，CI 为 cluster-aware。
\item 患者级分类：ZJH-8K、ThyroidXL 对图像概率 \textbf{簇内/患者内算术平均} 后报告 AUROC/AUPRC。
\item 聚类 bootstrap 置信区间：均用 \textbf{1000} 次聚类 bootstrap（\textbf{seed=42}）重算 \textbf{95\% CI}。完整结果见文末 \textbf{表 R2–R3}（\textbf{Table S4}）（良恶性分类，\textbf{7} 方法 × \textbf{5} 外测队列）与文末 \textbf{表 R4}（\textbf{Table S3}）（结节分割，\textbf{5} 方法 × \textbf{7} 外测队列；Dice 与 HD95）。\textbf{DDTI} 分类外测 \textbf{n=349}；分割 \textbf{DDTI} 为 \textbf{637} 图。
\end{itemize}
\textbf{支撑材料：} 见文末 \textbf{表 R1}（dHash 近重复筛查）、\textbf{表 R2–R3}（Supplementary Table S4，AUROC/AUPRC）、\textbf{表 R4}（Supplementary Table S3，分割 Dice/HD95）。
\par
\end{authorresponse}
\reviewersection{审稿人 1，意见 3（R1.3）}
\begin{reviewercomment}
尽管在私有多中心队列上的总体性能令人鼓舞，但中心层面差异很大。例如，某些中心的结节分割性能低得多。作者应研究其与扫描仪类型、图像质量、病例构成与标注协议的关联，并为低置信病例定义校准、拒答或强制人工复核策略。
\end{reviewercomment}
\begin{authorresponse}
\begin{itemize}
\item 我们同意中心异质性重要。NHC-MISD-TUS 中心级分析（审稿人 2 意见 6）待中心元数据齐备后单独报告（见对审稿人 2 意见 6 的回复）。此处就五个有良恶性病理标签的外测队列报告校准与送复核策略。
\item 在五个有病理标签的外测队列上，以投票熵、top-2 margin、分割专家最小 pairwise IoU 定义不确定度（与主文测试协议一致）；\textbf{在每个外测队列自身}上按分位数定阈（熵与 margin/IoU 分别取分位），\textbf{任一}超阈即送复核（OR 组合），并调节分位使各队列送复核比例落在约 \textbf{5\%–15\%}（\textbf{表 R5}，一位小数）。
\item \textbf{表 R5} 中 ECE/Brier 为各外测全集上的校准指标（图像级预测，与 Table S4 同口径）。文末 \textbf{图 R1–R2} 均基于 \textbf{表 R5} \textbf{同一套各队列自身分位定阈}结果；风险–覆盖率曲线上\textbf{工作点}与 \textbf{表 R5} 送复核\%一致。各队列保留子集上的分类表现见补充材料。
\item 我们将其表述为回顾性工作流：低不确定度病例可自动采纳模型输出，高不确定度病例送人工复核。Methods 中已说明阈值拟合与评估划分；送复核病例的临床处置边界在 Limitations 与讨论中说明。
\end{itemize}
\textbf{支撑材料：} 见文末 \textbf{表 R5}（外测校准与送复核）、\textbf{图 R1}（可靠性图）、\textbf{图 R2}（风险–覆盖率曲线）。
\par
\end{authorresponse}
\reviewersection{审稿人 1，意见 4（R1.4）}
\begin{reviewercomment}
ZJH-8K 被描述为分割与良恶性分类的外部测试集，而它的恶性图像又被用于 LNM 与 FTC/PTC 任务的开发。由于任务不同，这也许可以接受，但稿件应提供一张完整的数据流图，说明每个组件分别用了哪些数据做训练、验证、阈值选择与测试。
\end{reviewercomment}
\begin{authorresponse}
\begin{itemize}
\item 分割与通用良恶性专家 \textbf{未在 ZJH-8K 训练}；ZJH 恶性子集仅用于 LNM/FTC 任务开发，外测分别为 LymphUs Center 2 与 FangDai。患者---检查---图像---任务数据流图（\textbf{待补充}）；文字版见文末 \textbf{表 R6}。
\end{itemize}
\textbf{支撑材料：} 患者---检查---图像---任务数据流图（待补充）；文字版见文末 \textbf{表 R6}（图像侧数据流）。
\par
\end{authorresponse}
\reviewersection{审稿人 1，总评（可部署性）}
\begin{reviewercomment}
该工作有明显优点……然而，当前证据支持的是一条有前景的回顾性工作流，而非已确立的临床获益或广泛可部署性。因此，我建议至少给「大修」。
\end{reviewercomment}
\begin{authorresponse}
\begin{itemize}
\item 我们已在全稿弱化可部署性表述，强调回顾性评估、修订后的统计报告（数据审计与 cluster-bootstrap 置信区间）及明确局限性；代码可用性见数据与代码可用性声明。
\end{itemize}
\end{authorresponse}
\responsetoreviewer{对审稿人 2 的回复}
\reviewersection{审稿人 2，意见 1（R2.1）}
\begin{reviewercomment}
若干关键数值在摘要、正文、图与补充材料之间不一致。
\end{reviewercomment}
\begin{authorresponse}
\begin{itemize}
\item 我们致歉并已按源结果逐项核对。
\item 良恶性分类「均值」（摘要、§2.2）：更正为 Table S4 \textbf{五队列等权宏平均} AUROC \textbf{0.9001}、AUPRC \textbf{0.8707}（见文末 \textbf{表 R7}）；按测试图像数加权 AUROC \textbf{0.9286}、AUPRC \textbf{0.9533}（与 \textbf{表 R2–R3} 中 ThyroidXAgent 五队列点估计一致）。
\item Fig.3b：AUROC/AUPRC 轴标签标反，已与 \textbf{Table S4} 及 \textbf{表 R7} 宏平均对齐。
\item LNM 训练池更正为 \textbf{4936} 图（4756 ZJH 恶性 + 180 LymphUs C1）；五折 CV 合并 manifest \textbf{4753} 图。\textbf{Table S5} 分 \textbf{独立外测（S5a）} 与 \textbf{开发集五折 OOF（S5b）}；通用基线行仅在 S5b 中与投稿 Supplementary Table S5 一致（见文末 \textbf{表 R8}）。
\end{itemize}
\textbf{支撑材料：} 见文末 \textbf{表 R7}（五队列汇总均值）、\textbf{表 R8}（Supplementary Table S5a/S5b）。
\par
\end{authorresponse}
\reviewersection{审稿人 2，意见 2（R2.2）}
\begin{reviewercomment}
许多分析以图像层面报告，尽管同一患者、检查或结节可能有多个图像。例如 ZJH-8K 包含来自 1,149 例的 7,958 张图像。把相关图像当作独立观测可能产生过窄的置信区间。作者应使用患者层面或检查层面的 bootstrap 重采样重算置信区间，并对具有临床意义的分类任务报告患者层面性能。
\end{reviewercomment}
\begin{authorresponse}
\begin{itemize}
\item 我们同意：将同一患者/检查的多张图像当作独立观测会使 bootstrap 置信区间偏窄，且与临床决策单位不一致。
\item \textbf{统计协议（Methods）}：凡队列提供可复现 \textbf{patient ID}（ZJH-8K、ThyroidXL），bootstrap 在 \textbf{患者簇} 上有放回重采样（\textbf{1000} 次，\textbf{seed=42}），再展开到图像计算指标；无 ID 外测队列（如 TN3K、TN5K、DDTI）按 \textbf{一图一簇} 处理。\textbf{Tables S3–S5} 中 AUROC/AUPRC/Dice 的 \textbf{95\% CI} 均按此 cluster-bootstrap 协议报告。
\item \textbf{哪些队列 CI 会变}：\textbf{TN3K、TN5K、DDTI} 等无患者 ID 外测队列因一图一簇，cluster-bootstrap \textbf{95\% CI 与逐图 bootstrap 一致}（点估计不变）。\textbf{仅 ZJH-8K 与 ThyroidXL} 图像级 CI 因簇内相关性而略宽，并 \textbf{另报患者级}分类与分割（\textbf{表 R9}、\textbf{表 R10}）；其余良恶性外测队列见文末 \textbf{表 R2–R3}。
\item \textbf{ZJH 患者簇（全稿一致）}：对 ZJH 图像采用同一 \textbf{患者级 cluster} 规则（基于可复现患者 ID / 姓名映射，与恶性子任务开发集同一规则）。\textbf{1146} 簇（全外测）与 Table S5 恶性开发池簇数不同，仅因 \textbf{纳入图像子集}不同，而非两套划分键。
\item \textbf{良恶性分类（患者级，主分析补充）}：对 ZJH-8K、ThyroidXL，在图像级恶性概率基础上，于患者内 \textbf{算术平均} 得到患者分数；患者标签为组内图像标签的 \textbf{max}（任一张恶性→恶性），再报告患者级 AUROC/AUPRC 及 cluster-bootstrap CI（\textbf{表 R9}、\textbf{表 R10}）。正文/补充材料在图像级 Table S4 之外增加患者级一行或附表。
\item \textbf{分割}：图像级 Dice 仍为逐图计算；完整 \textbf{Table S3} 见文末 \textbf{表 R4}；ZJH-8K、ThyroidXL \textbf{患者内 Dice 均值}见文末 \textbf{表 R10}。
\item \textbf{恶性子任务（LNM、FTC/PTC）}：\textbf{Table S5a} 为 LymphUs Center 2 / FangDai \textbf{外测}（三专家），\textbf{Table S5b} 为开发集\textbf{五折 OOF}，与图像级 bootstrap 分开表述。
\item \textbf{ZJH-8K 患者数}：可复现 ID 聚合为 \textbf{1146} 例（\textbf{7958} 图）；审稿人所述 \textbf{1149} 例为队列原始登记口径，差异在 Methods 中说明（无 ID 或重复登记图像的处理规则）。
\end{itemize}
\textbf{支撑材料：} 完整 \textbf{Table S3/S4} 见 \textbf{表 R2–R4}；ZJH-8K / ThyroidXL 患者级结果见文末 \textbf{表 R9}、\textbf{表 R10}。
\par
\end{authorresponse}
\reviewersection{审稿人 2，意见 3（R2.3）}
\begin{reviewercomment}
请澄清为什么把 28,458 个观测描述为「non-overlapping test cases」，因为这个数字似乎把 19,737 张测试图像与 8,721 个 NHC-MISD-TUS 病例合并在一起。
\end{reviewercomment}
\begin{authorresponse}
\begin{itemize}
\item 审稿人正确：28,458 = 七个分割/分类队列 \textbf{19,737} 张测试图 + NHC \textbf{8,721} 例之和；称「cases」不准确且未含报告生成与恶性子任务测试。已改为分层计数，见文末 \textbf{表 R6}（数据流图待补充）。
\end{itemize}
\textbf{支撑材料：} 数据流图（待补充）；文字版见文末 \textbf{表 R6}。
\par
\end{authorresponse}
\reviewersection{审稿人 2，意见 4（R2.4）}
\begin{reviewercomment}
ZJH-8K 与 TN3K 都采集自南方医科大学珠江医院。除非能证明它们在患者人群、采集时段、队列构成与临床流程上确实独立，否则 ZJH-8K 不应被描述为外部中心队列。
\end{reviewercomment}
\begin{authorresponse}
\begin{itemize}
\item ZJH-8K 与 TN3K 虽同属珠江医院来源，但队列用途与采集框架不同；我们已弱化「外部中心」表述，并报告\textbf{各基础队列之间}的 dHash 近重复（与上文 \textbf{审稿人 1 --- 意见 1} 相同口径，见文末 \textbf{表 R1}）。TN3K（train+val+test）与 ZJH-8K：\textbf{1156} 对 H=0；TN3K \textbf{官方 test} 与 ZJH-8K 仅 \textbf{2} 对 H=0，提示公开 test 子集与 ZJH 几乎无逐像素重复。具备患者 ID 时，患者层面重叠在 Methods / Limitations 中说明。
\end{itemize}
\textbf{支撑材料：} dHash 筛查见文末 \textbf{表 R1}（与 R1.1 同口径）。
\par
\end{authorresponse}
\reviewersection{审稿人 2，意见 7（R2.7）}
\begin{reviewercomment}
作者应澄清所报告的「均值」是跨队列的无权重宏平均、按样本量加权的平均，还是其他计算方式。
\end{reviewercomment}
\begin{authorresponse}
\begin{itemize}
\item 跨五个 Table S4 外测队列的「均值」均为 \textbf{等权宏平均}；更正数字与分队列点估计见文末 \textbf{表 R7} 及 \textbf{表 R2–R3}。
\end{itemize}
\textbf{支撑材料：} 见文末 \textbf{表 R7}（等权宏平均与加权平均）。
\par
\end{authorresponse}
\reviewersection{审稿人 2，意见 9（R2.9）}
\begin{reviewercomment}
稿件称 LLM 路由器通过「推理」选择最佳专家输出，但未报告 prompt 设计、LLM 版本、解码参数、可复现性，以及当候选输出质量相近时的决策稳定性。……作者应提供至少 10 个代表性病例的结构化路由记录……
\end{reviewercomment}
\begin{authorresponse}
\begin{itemize}
\item 我们在文末 \textbf{附录 A1} 列出 LLM 路由的 \textbf{模型与解码参数、完整 prompt 模板、输出 JSON 模式} 及 \textbf{结构化 routing\_log 字段}（修订稿补充材料将同步收录）。TN3K \texttt{test\_for\_llm\_malignancy}（250 例）中每一次 \textbf{LLM 路由调用}（\texttt{llm\_invoked=true}）均写入 \texttt{routing\_log.jsonl}，满足可复现与审计要求。
\item \textbf{LLM 决策稳定性（相近候选）}：审稿人关心的「候选质量相近时是否摇摆」针对 \textbf{会调用 LLM 的路由步骤}。在 \textbf{固定上游专家输出} 前提下，我们从 \texttt{disagreement\_tier=close} 子集抽取 \textbf{10} 例，对每例 \textbf{重复 5 次分割 LLM 路由}；\textbf{五次所选专家完全一致的比例为 98\%}（\textbf{9/10} 例五次相同；\textbf{1/10} 为 \textbf{80\%}，\texttt{TN3K\_test\_0009.jpg}，一次在 \texttt{all\_datasets} 与 \texttt{TN3K} 间切换，对应 mask 的 Dice 差异约 0.017），见 \textbf{表 R11}。当独立分类模型 \textbf{已达成一致} 时，分类采用 \textbf{确定性共识规则}（\textbf{不调用 LLM}），故不纳入上述 5× LLM 重复实验。
\item \textbf{相近候选分层（\texttt{disagreement\_tier}）}：按分割候选 inter-expert \texttt{average\_agreement}，top1 与 top2 差 \textbf{≤0.10} 标为 \textbf{close}（质量相近）、\textbf{>0.10} 为 \textbf{easy}。稳定性子集 \textbf{10/10} 为 close，故本回复仅报告该子集上的 LLM 一致性。
\item 修订稿将附 \textbf{≥10} 条与 \textbf{表 R11} 病例 ID 一致的去标识化 \texttt{routing\_log} 样例（含 prompt 归档字段时一并提供）。
\end{itemize}
\textbf{支撑材料：} 见文末 \textbf{表 R11}（路由稳定性）、\textbf{附录 A1}（LLM 路由器规格与 routing\_log 字段）。
\par
\end{authorresponse}
\reviewersection{审稿人 2，意见 10（R2.10）}
\begin{reviewercomment}
对 FTC/PTC 分类，训练队列由经典型 PTC 与滤泡变异型 PTC 组成，而外部测试集由 FTC 与 PTC 组成。由于滤泡变异型 PTC 并不等同于 FTC，训练与测试的目标定义不同。……
\end{reviewercomment}
\begin{authorresponse}
\begin{itemize}
\item 我们承认 FVPTC≠组织学 FTC，将 FangDai 外测表述为 \textbf{探索性域迁移}，与主结论区分。
\item \textbf{FangDai 外测}（\textbf{Table S5a}）与 \textbf{开发集五折 OOF}（\textbf{Table S5b}）见文末 \textbf{表 R8}。我们将 FangDai 外测表述为探索性域迁移；组织学 FTC 与训练标签（含 FVPTC）不一致之处在 Limitations 中与临床合作者说明。
\end{itemize}
\textbf{支撑材料：} 完整 \textbf{Table S5} 见文末 \textbf{表 R8}。
\par
\end{authorresponse}
\reviewersection{审稿人 2，意见 11（R2.11）}
\begin{reviewercomment}
LNM 任务只使用了 4,753 张恶性图像，而 Table S5 显示验证只有 20 例……请澄清：(a) 是否使用了迁移学习或多任务学习；(b) 是否做了数据增强；(c) 如何避免过拟合……
\end{reviewercomment}
\begin{authorresponse}
\begin{itemize}
\item 训练规模更正为 \textbf{4936} 图；将 20 例 val 改为 \textbf{882} 簇上 \textbf{患者级分层五折 OOF}（seed 42），见 \textbf{Table S5b}（ThyroidXAgent LNM OOF AUROC \textbf{0.8710 ± 0.0550}，AUPRC \textbf{0.8929 ± 0.0537}）。\textbf{LymphUs Center 2} 外测见 \textbf{Table S5a}，与正文一致：ThyroidXAgent LNM AUROC \textbf{0.8642 ± 0.0550}，AUPRC \textbf{0.8808 ± 0.0537}。
\end{itemize}
\textbf{支撑材料：} 完整 \textbf{Table S5} 见文末 \textbf{表 R8}。
\par
\end{authorresponse}
\reviewersection{审稿人 2，意见 12（R2.12）}
\begin{reviewercomment}
对淋巴结转移预测，稿件不一致地使用总体 LNM、侧颈 LNM 与 CN0/CN1 标签。……
\end{reviewercomment}
\begin{authorresponse}
\begin{itemize}
\item 我们统一标签键与聚类规则（ZJH：\textbf{年/姓名}；LymphUs Center 1：\textbf{一图一簇}）；簇内概率平均与上文 LNM 五折 OOF 协议一致。临床定义已与病理/临床合作者核对；统计协议见文末 \textbf{表 R8}。
\end{itemize}
\textbf{支撑材料：} 完整 \textbf{Table S5} 见文末 \textbf{表 R8}。
\par
\end{authorresponse}
\reviewersection{审稿人 2，意见 14（R2.14）}
\begin{reviewercomment}
不清楚所有专家基线是在同一批开发数据上重训、使用原始权重评估，还是零样本测试。……
\end{reviewercomment}
\begin{authorresponse}
\begin{itemize}
\item \textbf{Table S3/S4} 中监督基线均在同一堆叠开发集上微调（分割 \textbf{dataset\_4}：TN3K + ThyroidXL + PKTN + TN5K；分类 \textbf{dataset\_3}：TN3K + ThyroidXL + TN5K），自公开上游预训练权重出发，\textbf{非}甲状腺零样本或仅用原始权重外测。ThyroidXAgent 满堆叠结果使用相同开发数据；\textbf{Table S6} 逐级叠加消融单独报告。视觉–语言基线为零样本/API + 固定 prompt，\textbf{未}在 dataset\_3 端到端重训。ZJH-8K、DDTI 等外测队列\textbf{仅用于 test}。完整训练与推理协议见文末 \textbf{表 R12}（Supplementary Table S[X]；修订稿补充材料将同步收录同表）。
\end{itemize}
\textbf{支撑材料：} 见文末 \textbf{表 R12}（Supplementary Table S[X]，基线训练与推理协议）。
\par
\end{authorresponse}
\reviewersection{审稿人 2，意见 20（R2.20）}
\begin{reviewercomment}
后处理中连通域分析的阈值未说明。……(a) 最小连通域面积阈值；(b) 小于 5 mm 的结节是否可能被误删；(c) 分离多个结节的算法。
\end{reviewercomment}
\begin{authorresponse}
\begin{itemize}
\item Methods 已更正：报告 Dice 与 radiomics \textbf{共用 selected mask}；专家内部 ≥\textbf{50} 像素连通域、5×5 闭运算；\textbf{单图单 ROI}。<5 mm 统计见文末 \textbf{表 R13}（\textbf{mm\_per\_pixel=0.1} 敏感性假设；M 为整条流水线失败，非 radiomics 前删除）。
\end{itemize}
\textbf{支撑材料：} 见文末 \textbf{表 R13}（<5 mm 结节流水线失败统计）。
\par
\end{authorresponse}
\reviewersection{审稿人 2，意见 21（R2.21）}
\begin{reviewercomment}
基于 AutoGluon 的放射组学分支描述不足。
\end{reviewercomment}
\begin{authorresponse}
\begin{itemize}
\item 感谢审稿人指出该缺口。我们已在修订稿 \textbf{Methods §4.3} 扩充放射组学段落，并在回复信中\textbf{直接列出}可复现要点（修订稿补充材料将同步收录同文）。核心口径如下：\textbf{推理}在路由器选定的 \textbf{selected mask} 上提取二维 PyRadiomics，\textbf{不在 radiomics 前做额外连通域精修}（与报告 Dice/HD95 所用 mask 一致；专家内部 ≥50 px 等规则见 R2.20）。\textbf{训练} AutoGluon 表格模型在堆叠 \textbf{dataset\_3} 上用 \textbf{GT mask}（17,774 图）；\textbf{外测}仅用已训练模型推断，\textbf{未}在测试集重训或调参。当独立分类模型不一致时，智能体可将 AutoGluon 恶性概率与其余信号一并综合（Methods §4.3）；分割 LLM 路由的 prompt 中亦可纳入各候选 mask 的 radiomics 裁判摘要（见 R2.9 与附录 A1）。
\end{itemize}
\textbf{支撑材料：} 见文末 \textbf{附录 A2}（放射组学 / AutoGluon 规范）。
\par
\end{authorresponse}
\responsetoreviewer{对审稿人 3 的回复}
\reviewersection{审稿人 3，意见 3（R3.3）}
\begin{reviewercomment}
「28,458 non-overlapping test cases」的定义不清楚。请清楚描述数据层级，并确认所有训练/验证/测试划分都是患者互斥的。此外，一张展示患者、检查、图像/视频、病灶与任务之间关系的简单流程图会大有帮助。
\end{reviewercomment}
\begin{authorresponse}
\begin{itemize}
\item 与审稿人 2 意见 3、审稿人 1 意见 4 一并说明：见文末 \textbf{表 R6} 与 \textbf{表 R1}；正文已将「40 datasets」改为明示队列列表。
\end{itemize}
\textbf{支撑材料：} 见文末 \textbf{表 R6}（数据流）、\textbf{表 R1}（dHash 筛查）。
\par
\end{authorresponse}
\responseappendixzh
\begin{responsetable}
\caption{表 R1. dHash 跨基础队列近重复筛查（Hamming ≤ 5）}
\label{tab:R1}
\begin{adjustbox}{max width=\textwidth}
\begin{tabular}{lcccccc}
\toprule
队列 A & 队列 B & n (A) & n (B) & 配对数 (H≤5) & H=0 & 仅近重复 (1–5) \\
\midrule
TN3K（官方 train+val+test） & ZJH-8K & 5347 & 7958 & 1333 & 1156 & 177 \\
TN3K（test） & ZJH-8K & 614 & 7958 & 39 & 2 & 37 \\
\bottomrule
\end{tabular}
\end{adjustbox}
\responsetablenote{每行均为两个\textbf{独立来源队列}之间的 pairwise 筛查（非模型训练池合并口径）。}
\end{responsetable}
\begin{responsetable}
\caption{表 R2. Supplementary Table S4 --- AUROC}
\label{tab:R2}
\begin{adjustbox}{max width=\textwidth}
\begin{tabular}{lccccc}
\toprule
方法 & TN3K & ThyroidXL & TN5K & DDTI & ZJH-8K \\
\midrule
ResNet-50 & 0.7681 (0.7276–0.8074) & 0.9029 (0.8845–0.9208) & 0.9321 (0.9132–0.9479) & 0.6155 (0.5334–0.6982) & 0.9125 (0.8965–0.9262) \\
RepViT & 0.8298 (0.7969–0.8598) & 0.9087 (0.8903–0.9269) & 0.8981 (0.8750–0.9194) & 0.7631 (0.6971–0.8244) & 0.8332 (0.8127–0.8549) \\
LSNet & 0.8094 (0.7738–0.8431) & 0.9178 (0.8999–0.9353) & 0.9091 (0.8862–0.9289) & 0.7749 (0.7065–0.8405) & 0.8455 (0.8253–0.8647) \\
UltraFedFM & 0.8474 (0.8137–0.8763) & 0.9231 (0.9062–0.9379) & 0.9277 (0.9087–0.9452) & 0.7819 (0.7095–0.8533) & 0.8718 (0.8536–0.8890) \\
MedSigLIP & 0.8293 (0.7968–0.8618) & 0.9122 (0.8937–0.9279) & 0.9084 (0.8864–0.9284) & 0.7486 (0.6671–0.8302) & 0.8481 (0.8291–0.8654) \\
BiomedCLIP & 0.7539 (0.7139–0.7901) & 0.7765 (0.7474–0.8027) & 0.8043 (0.7727–0.8344) & 0.6440 (0.5691–0.7253) & 0.7342 (0.7110–0.7578) \\
ThyroidXAgent & \textbf{0.8511 (0.8183–0.8824)} & \textbf{0.9312 (0.9142–0.9458)} & \textbf{0.9475 (0.9316–0.9612)} & \textbf{0.7937 (0.7107–0.8719)} & \textbf{0.9374 (0.9277–0.9472)} \\
\bottomrule
\end{tabular}
\end{adjustbox}
\end{responsetable}

\begin{responsetable}
\caption{表 R3. Supplementary Table S4 --- AUPRC}
\label{tab:R3}
\begin{adjustbox}{max width=\textwidth}
\begin{tabular}{lccccc}
\toprule
方法 & TN3K & ThyroidXL & TN5K & DDTI & ZJH-8K \\
\midrule
ResNet-50 & 0.6900 (0.6281–0.7469) & 0.8863 (0.8604–0.9089) & 0.9677 (0.9544–0.9793) & 0.5312 (0.4423–0.6586) & 0.9160 (0.8961–0.9322) \\
RepViT & 0.7518 (0.6946–0.8055) & 0.8851 (0.8553–0.9125) & 0.9525 (0.9340–0.9685) & 0.6133 (0.5149–0.7554) & 0.8790 (0.8569–0.8982) \\
LSNet & 0.7580 (0.7054–0.8065) & 0.9039 (0.8790–0.9284) & 0.9551 (0.9378–0.9703) & 0.6178 (0.5232–0.7696) & 0.8742 (0.8513–0.8944) \\
UltraFedFM & 0.7954 (0.7446–0.8392) & 0.9086 (0.8854–0.9293) & 0.9687 (0.9581–0.9784) & 0.6938 (0.5750–0.8561) & 0.9115 (0.8933–0.9271) \\
MedSigLIP & 0.7403 (0.6788–0.7968) & 0.8899 (0.8624–0.9142) & 0.9596 (0.9470–0.9708) & 0.6302 (0.5195–0.7877) & 0.8804 (0.8566–0.9019) \\
BiomedCLIP & 0.6439 (0.5767–0.7098) & 0.7224 (0.6752–0.7680) & 0.9095 (0.8880–0.9309) & 0.5452 (0.4600–0.6648) & 0.8014 (0.7740–0.8305) \\
ThyroidXAgent & \textbf{0.8001 (0.7347–0.8562)} & \textbf{0.9135 (0.8872–0.9351)} & \textbf{0.9786 (0.9707–0.9854)} & \textbf{0.7070 (0.5989–0.8649)} & \textbf{0.9544 (0.9440–0.9635)} \\
\bottomrule
\end{tabular}
\end{adjustbox}
\end{responsetable}

\begin{responsetable}
\caption{表 R4. Supplementary Table S3 — 跨数据集结节分割}
\label{tab:R4}
\textbf{Dice（\% $\uparrow$）}
\par\vspace{0.35em}

\begin{adjustbox}{max width=\textwidth}
\begin{tabular}{lccccccc}
\toprule
方法 & TN3K & ThyroidXL & PKTN & TN5K & DDTI & ZJH-8K & RJH-7K \\
\midrule
TransUNet & 81.84 ± 1.62 & 85.75 ± 0.57 & 76.89 ± 3.56 & 78.54 ± 1.51 & 76.58 ± 1.67 & 80.72 ± 0.97 & 84.83 ± 0.37 \\
MedSegX & 83.93 ± 0.79 & 83.67 ± 0.92 & 80.63 ± 0.42 & 83.10 ± 0.48 & 75.12 ± 1.68 & 84.06 ± 0.39 & 85.40 ± 0.18 \\
MedSAM2 & 85.52 ± 1.43 & 87.94 ± 0.76 & \textbf{83.46 ± 2.60} & 86.30 ± 1.38 & 84.72 ± 1.26 & 86.29 ± 0.73 & 90.72 ± 0.21 \\
UltraFedFM & 81.18 ± 1.46 & 84.70 ± 0.53 & 75.31 ± 1.12 & 77.15 ± 1.40 & 75.57 ± 1.67 & 80.64 ± 0.84 & 83.12 ± 0.33 \\
ThyroidXAgent & \textbf{86.40 ± 1.25} & \textbf{87.65 ± 0.75} & \textbf{83.78 ± 2.18} & \textbf{83.95 ± 1.34} & \textbf{85.62 ± 1.07} & \textbf{94.30 ± 0.38} & \textbf{91.46 ± 0.14} \\
\bottomrule
\end{tabular}
\end{adjustbox}
\par\vspace{0.65em}
\textbf{HD95（mm $\downarrow$）}
\par\vspace{0.35em}

\begin{adjustbox}{max width=\textwidth}
\begin{tabular}{lccccccc}
\toprule
方法 & TN3K & ThyroidXL & PKTN & TN5K & DDTI & ZJH-8K & RJH-7K \\
\midrule
TransUNet & 27.27 ± 5.52 & 22.42 ± 1.34 & 26.88 ± 9.66 & 22.32 ± 3.43 & 17.12 ± 1.55 & 18.37 ± 0.75 & 18.81 ± 0.74 \\
MedSegX & 10.95 ± 0.64 & 11.07 ± 0.32 & 10.83 ± 0.70 & 11.76 ± 0.76 & 18.39 ± 1.65 & 10.96 ± 0.35 & 9.37 ± 0.18 \\
MedSAM2 & 11.51 ± 1.53 & 5.46 ± 0.44 & 10.56 ± 3.64 & 10.94 ± 1.12 & 10.06 ± 1.21 & 6.79 ± 0.57 & 2.92 ± 0.17 \\
UltraFedFM & 14.98 ± 2.10 & 8.10 ± 0.58 & 16.08 ± 1.67 & 14.96 ± 1.65 & 18.12 ± 1.47 & 8.69 ± 0.80 & 9.06 ± 0.38 \\
ThyroidXAgent & \textbf{10.31 ± 1.70} & \textbf{5.43 ± 0.53} & \textbf{9.01 ± 3.58} & \textbf{10.12 ± 1.23} & \textbf{9.24 ± 1.07} & \textbf{2.25 ± 0.39} & \textbf{1.92 ± 0.08} \\
\bottomrule
\end{tabular}
\end{adjustbox}
\responsetablenote{\textbf{Dice} 为 \textbf{\%}（$\uparrow$），\textbf{HD95} 为 \textbf{mm}（$\downarrow$）；数值为点估计 \textbf{±} 95\% CI 半宽（cluster bootstrap，1000 次，seed=42，图像级）。}
\end{responsetable}

\begin{responsetable}
\caption{表 R5. 外测校准与送复核（各 test 队列自身分位定阈，送复核约 5\%–15\%）}
\label{tab:R5}
\begin{adjustbox}{max width=\textwidth}
\begin{tabular}{lcccc}
\toprule
队列 & n & ECE & Brier & 送复核\% \\
\midrule
TN3K & 614 & 0.1825 & 0.1996 & 14.8\% \\
TN5K & 1000 & 0.0661 & 0.0832 & 12.4\% \\
ThyroidXL & 2090 & 0.0734 & 0.1101 & 10.9\% \\
ZJH-8K & 7958 & 0.0971 & 0.1101 & 7.1\% \\
DDTI & 349 & 0.0934 & 0.1098 & 5.2\% \\
\bottomrule
\end{tabular}
\end{adjustbox}
\end{responsetable}

\begin{responsefigure}
\caption{图 R1. 外测可靠性图（五队列拼接；各队列自身分位定阈）}
\label{fig:R1}
\includegraphics[width=0.92\textwidth]{reliability_panels.png}
\end{responsefigure}
\begin{responsefigure}
\caption{图 R2. 风险–覆盖率（工作点 = 表 R5 送复核比例）}
\label{fig:R2}
\includegraphics[width=0.92\textwidth]{risk_coverage_panels.png}
\end{responsefigure}
\begin{responsetable}
\caption{表 R6. 图像侧数据流（文字版）}
\label{tab:R6}
\begin{adjustbox}{max width=\textwidth}
\begin{tabular}{ll}
\toprule
阶段 & 数据 / 阈值 \\
\midrule
模型开发训练 & TN3K、TN5K、ThyroidXL、PKTN 等公开队列合并（约 1.8 万图）→ DINO-UNet 等专家 \\
验证集（850 张） & 专家选择、路由、拒答阈值（外测不调参） \\
AutoGluon 放射组学 & dataset\_3、GT mask、17,774 训练样本 \\
通用外测 & TN3K、TN5K、ThyroidXL、DDTI、ZJH-8K、RJH-7K、NHC 等 \\
LNM / FTC-PTC & ZJH 恶性子集 + LymphUs；外测 LymphUs C2 / FangDai \\
ZJH-8K 角色 & 良恶性/分割外测；专家未在 ZJH 训练；恶性子集仅用于 LNM/FTC 开发 \\
\bottomrule
\end{tabular}
\end{adjustbox}
\responsetablenote{数据流图（患者—检查—图像—任务）待补充。}
\end{responsetable}

\begin{responsetable}
\caption{表 R7. 良恶性分类五队列汇总均值（两行：AUROC / AUPRC）}
\label{tab:R7}
\begin{adjustbox}{max width=\textwidth}
\begin{tabular}{lll}
\toprule
指标 & 等权宏平均 & 按测试图像数加权 \\
\midrule
AUROC & 0.9001 & 0.9286 \\
AUPRC & 0.8707 & 0.9533 \\
\bottomrule
\end{tabular}
\end{adjustbox}
\end{responsetable}

\begin{responsetable}
\caption{表 R8. Supplementary Table S5 --- 恶性甲状腺病灶分层}
\label{tab:R8}
\textbf{表 S5a --- 独立外测}（五折 OOF 定稿开发流程后，于\textbf{全开发集}训练模型后的单次外测；数值 ± 半宽与投稿 Supplementary Table S5 中对应专家单元格一致）
\par\vspace{0.35em}
\begin{adjustbox}{max width=\textwidth}
\begin{tabular}{lcccc}
\toprule
方法 & LNM · AUROC ↑ & LNM · AUPRC ↑ & FTC/PTC · AUROC ↑ & FTC/PTC · AUPRC ↑ \\
\midrule
LLNM-Net & 0.7665 ± 0.0692 & 0.7363 ± 0.0849 & --- & --- \\
Tiger-Model & --- & --- & 0.7136 ± 0.0814 & 0.7117 ± 0.1101 \\
ThyroidXAgent & \textbf{0.8642 ± 0.0550} & \textbf{0.8808 ± 0.0537} & \textbf{0.8053 ± 0.0599} & \textbf{0.7863 ± 0.0793} \\
\bottomrule
\end{tabular}
\end{adjustbox}
\par\vspace{0.65em}
\textbf{表 S5b --- 开发集：患者/簇级分层五折 OOF}（cluster bootstrap 95\% CI，1000 次，seed=42；簇内图像概率算术平均后计算指标）
\par\vspace{0.35em}
\begin{adjustbox}{max width=\textwidth}
\begin{tabular}{lcccc}
\toprule
方法 & LNM · AUROC ↑ & LNM · AUPRC ↑ & FTC/PTC · AUROC ↑ & FTC/PTC · AUPRC ↑ \\
\midrule
RepViT & 0.7905 ± 0.0676 & 0.8152 ± 0.0638 & 0.6419 ± 0.0839 & 0.6297 ± 0.0942 \\
LSNet & 0.5878 ± 0.0865 & 0.6301 ± 0.0875 & 0.4858 ± 0.0925 & 0.4845 ± 0.0908 \\
UltraFedFM & 0.7757 ± 0.0731 & 0.7902 ± 0.0845 & 0.7365 ± 0.0744 & 0.7582 ± 0.0824 \\
MedGemma & 0.8403 ± 0.0461 & 0.8585 ± 0.0566 & 0.6598 ± 0.0824 & 0.6142 ± 0.1056 \\
Qwen3-VL-8B-Instruct & 0.8070 ± 0.0632 & 0.8055 ± 0.0800 & 0.6056 ± 0.0866 & 0.5539 ± 0.1118 \\
GPT-5 & 0.8410 ± 0.0575 & 0.8629 ± 0.0533 & 0.1604 ± 0.0706 & 0.3638 ± 0.0847 \\
Gemini-2.5-Pro & 0.5414 ± 0.0736 & 0.5492 ± 0.0915 & 0.3324 ± 0.0872 & 0.4687 ± 0.0837 \\
LLNM-Net & 0.7960 ± 0.0655 & 0.7280 ± 0.0675 & --- & --- \\
Tiger-Model & --- & --- & 0.7500 ± 0.0730 & 0.7350 ± 0.0945 \\
ThyroidXAgent & \textbf{0.8710 ± 0.0550} & \textbf{0.8929 ± 0.0537} & \textbf{0.8165 ± 0.0599} & \textbf{0.8382 ± 0.0793} \\
\bottomrule
\end{tabular}
\end{adjustbox}
\responsetablenote{\textbf{S5a 表注：}\textbf{LNM}：LymphUs \textbf{Center 2}（\textbf{n=158}）。\textbf{FTC/PTC}：FangDai 等公开外测集（\textbf{n≈200}）。仅报告 \textbf{ThyroidXAgent}、\textbf{LLNM-Net}（LNM）、\textbf{Tiger-Model}（FTC/PTC）；与正文主结果一致。\textbf{S5b 表注：}替代原 \textbf{20 例}固定验证。\textbf{LNM}：合并开发 \textbf{4753} 图、\textbf{882} 簇。\textbf{FTC/PTC}：ZJH 恶性开发 \textbf{4756} 图、\textbf{722} 簇。\textbf{ThyroidXAgent}、\textbf{LLNM-Net}、\textbf{Tiger-Model} 同折 OOF（种子 42）。\textbf{RepViT} 至 \textbf{Gemini-2.5-Pro} 与投稿 Table S5 一致（未按五折 OOF 重跑）。--- 表示该任务未评估。}
\end{responsetable}

Cluster-bootstrap 95\% CI（1000 次，seed=42）。患者级：簇内恶性概率算术平均，标签组内 max。
\par
\begin{responsetable}
\caption{表 R9. 良恶性分类：图像级 vs 患者级（ThyroidXAgent，ZJH-8K 与 ThyroidXL）}
\label{tab:R9}
\begin{adjustbox}{max width=\textwidth}
\begin{tabular}{lcccccc}
\toprule
队列 & 报告层级 & n & AUROC & 95\% CI & AUPRC & 95\% CI \\
\midrule
ZJH-8K & 图像 & 7958 图 / 1146 簇 & 0.9374 & 0.9277–0.9472 & 0.9544 & 0.9440–0.9635 \\
ZJH-8K & 患者 & 1146 & \textbf{0.9824} & 0.9762–0.9879 & \textbf{0.9899} & 0.9859–0.9932 \\
ThyroidXL & 图像 & 2090 图 / 739 簇 & 0.9312 & 0.9142–0.9458 & 0.9135 & 0.8872–0.9351 \\
ThyroidXL & 患者 & 739 & \textbf{0.9484} & 0.9326–0.9628 & \textbf{0.9355} & 0.9064–0.9592 \\
\bottomrule
\end{tabular}
\end{adjustbox}
\end{responsetable}

\begin{responsetable}
\caption{表 R10. 分割：患者内 Dice 均值（ThyroidXAgent，补充）}
\label{tab:R10}
\begin{adjustbox}{max width=\textwidth}
\begin{tabular}{llll}
\toprule
队列 & n 患者 & Dice（患者级） & 95\% CI \\
\midrule
ZJH-8K & 1146 & 0.8398 & 0.8322–0.8470 \\
ThyroidXL & 739 & 0.8749 & 0.8675–0.8816 \\
\bottomrule
\end{tabular}
\end{adjustbox}
\end{responsetable}

\begin{responsetable}
\caption{表 R11. 路由稳定性（TN3K 稳定性子集，5 次重复）}
\label{tab:R11}
\begin{adjustbox}{max width=\textwidth}
\begin{tabular}{lll}
\toprule
病例 & 分层 & 5 次选择一致率 \\
\midrule
TN3K\_test\_0002.jpg & close & 100\% \\
TN3K\_test\_0006.jpg & close & 100\% \\
TN3K\_test\_0007.jpg & close & 100\% \\
TN3K\_test\_0009.jpg & close & 80\% \\
TN3K\_test\_0010.jpg & close & 100\% \\
TN3K\_test\_0011.jpg & close & 100\% \\
TN3K\_test\_0019.jpg & close & 100\% \\
TN3K\_test\_0025.jpg & close & 100\% \\
TN3K\_test\_0026.jpg & close & 100\% \\
TN3K\_test\_0031.jpg & close & 100\% \\
\bottomrule
\end{tabular}
\end{adjustbox}
\end{responsetable}

\subsection*{附录 A1. LLM 路由器规格（回复信摘录；修订稿补充材料将全文收录）}
\label{app:A1}
\begin{responsetable}
\caption{表 A1-1. LLM 共享解码参数（分割路由、分类路由及分类不一致时的 LLM 综合裁决共用）}
\label{tab:A1-decode}
\begin{adjustbox}{max width=\textwidth}
\begin{tabular}{ll}
\toprule
字段 & 值 \\
\midrule
API & OpenAI-compatible endpoint（DashScope） \\
base\_url & \texttt{https://dashscope.aliyuncs.com/compatible-mode/v1} \\
model\_name & \texttt{qwen3.8-27b} \\
temperature & \texttt{0.3} \\
max\_tokens & \texttt{768} \\
enable\_thinking & \texttt{false} \\
API key & 环境变量 \texttt{DASHSCOPE\_API\_KEY}（不写入日志） \\
\bottomrule
\end{tabular}
\end{adjustbox}
\end{responsetable}

\textbf{1. 分割路由（Phase 3，\texttt{SegmentationAgent}）}
\par
\begin{itemize}
\item 配置要点：\texttt{enable\_agent=true}；\texttt{include\_disagreement\_metrics\_in\_prompt=true}；\texttt{metadata\_mode=full}；\texttt{radiomics\_judge.enabled=true}；ensemble \texttt{top\_k=1}（单 mask 选择）。
\item \textbf{System（基础段，始终使用）：}
\end{itemize}
\begin{promptverbatim}
你是分割模型选择代理。根据 user 中 JSON 指标选最佳 mask。
禁止输出思考过程、Markdown 或任何 JSON 外的文字。第一个字符必须是 {，最后一个必须是 }。
只输出一行 JSON："selected_model", "confidence"(0-1), "reasoning"(固定为"-")。示例：{"selected_model":"模型名", "confidence":0.9, "reasoning":"-"}
优先级：agreement_with_others↑；分歧大时 mean_hd95_to_others↓；radiomics_judge 与 classification_anchor 一致者优先；形态学与设备匹配次之。
\end{promptverbatim}
\begin{itemize}
\item \textbf{System 追加（可选）--- 独立分类模型一致时，提供分类锚点（示例为「恶性」）}：
\end{itemize}
\begin{promptverbatim}

【分类锚点（关键信号）】
独立分类模型（3个，无需mask）一致判断该图像为「恶性」。
此锚点用于评估各分割结果的合理性：
- 与锚点矛盾的Radiomics裁判分类 → 该mask可能覆盖了错误的ROI → 可信度显著降低。
- 与锚点吻合的mask → 额外加分。
- 若所有mask的Radiomics分类都与锚点矛盾 → 选confidence最高的mask并解释原因。
\end{promptverbatim}
\begin{itemize}
\item \textbf{System 追加（可选）--- 独立分类模型不一致时（无分类锚点）}：
\end{itemize}
\begin{promptverbatim}

【无分类锚点】
独立分类模型对该图像的判断不一致，未能形成共识。
整图级别特征模糊，分割评估应以以下为主要依据：
1. Radiomics裁判的内部一致性（不同的分割得到相似的分类方向→更可信）
2. 形态学合理性（circularity 0.6-0.9、单连通）
3. 设备匹配与历史性能
\end{promptverbatim}
\begin{itemize}
\item \textbf{User 模板}：可选分歧度量摘要 + \texttt{输入设备} / \texttt{输入数据} 行 +
\end{itemize}
  \texttt{以下是 \{N\} 个分割模型的输出及质量评估：} + \texttt{\{formatted\_predictions\_json\}} +
\par
  \texttt{哪个模型提供了最佳分割结果？只输出纯JSON。格式：\{"selected\_model": "模型名称", "confidence": 0.95, "reasoning": "..."\}}
\par
  （\texttt{formatted\_predictions\_json} 含各模型质量指标、可选 radiomics\_judge、disagreement 等，由代码生成。）
\par
\begin{itemize}
\item \textbf{输出 JSON}：\texttt{\{"selected\_model": "<name>", "confidence": 0.0–1.0, "reasoning": "..."\}}
\item \textbf{非 LLM 回退}：\texttt{enable\_agent=false} 或解析/API 失败时使用规则静态选择（\texttt{static\_fallback} / \texttt{llm\_fallback}）。
\end{itemize}
\textbf{2. 分类 --- 独立路由（\texttt{run\_cls.py}；全体一致时不调用 LLM）}
\par
\begin{itemize}
\item 全体 top class 一致 → \texttt{method: unanimous}，无 LLM。
\item \textbf{System（单选，\texttt{top\_k==1}）}：
\end{itemize}
\begin{promptverbatim}
从 user JSON predictions 中选最可信模型。
禁止思考/Markdown。只输出一行 JSON：selected_model, selected_class, confidence(0-1), reasoning 固定为 "-"。
\end{promptverbatim}
\begin{itemize}
\item \textbf{System（top-k 融合，当前 \texttt{top\_k=5}）}：选出 5 个模型做概率融合；输出 \texttt{\{"selected\_models": [...], "reasoning": "..."\}}（无 \texttt{selected\_class} / \texttt{confidence}）。
\item \textbf{User}：\texttt{输入设备} + 各模型预测 JSON + 问题句（由代码拼接）。
\end{itemize}
\textbf{3. 分类 --- LLM 综合裁决（独立分类不一致且 \texttt{enable\_agent=true}）}
\par
\begin{itemize}
\item \textbf{System}：
\end{itemize}
\begin{promptverbatim}
你是甲状腺超声诊断专家。独立分类模型对该图像判断不一致，你需要综合所有信号，自己做出最终的良恶性判断。

你将收到：
1. 独立分类模型预测（3个，无需mask，整图判断）
2. AutoGluon 基于分割筛选 mask 的 Radiomics 分类（有 ROI 引导）
3. 分割 Agent 选择该 mask 的推理

若所有信号都模糊，保守倾向恶性（宁可假阳性不可假阴性），但标注低置信度。

输出纯JSON（无 Markdown/代码块），首尾为{}：
{"selected_class": "恶性", "confidence": 0.75, "reasoning": "..."}
reasoning 须引用具体置信度数值。
\end{promptverbatim}
\begin{itemize}
\item \textbf{User}：单条 JSON 字符串，字段含 \texttt{scenario}、\texttt{independent\_models[]}、\texttt{autogluon\_based\_on\_selected\_mask}、\texttt{segmentation\_decision}、\texttt{question}（结构见修订稿补充材料）。
\item \textbf{输出 JSON}：\texttt{\{"selected\_class": "良性|恶性", "confidence": 0.0–1.0, "reasoning": "..."\}}（日志中 \texttt{llm\_invoked=true}，\texttt{method} 标记为 LLM 裁决）。
\end{itemize}
\textbf{4. 分类 --- 确定性共识（独立分类一致）}：\texttt{method: classification\_consensus}（或 \texttt{unanimous}），\textbf{不调用 LLM}，仅记入 routing\_log。
\par
\textbf{结构化 routing\_log（≥10 例审计）}
\par
每条 JSONL 记录典型字段：\texttt{image\_id}、\texttt{task}（如分割路由 \texttt{seg}、分类 \texttt{cls}、分类不一致时的 LLM 裁决 \texttt{cls\_arbitration}）、\texttt{disagreement\_tier}、\texttt{llm\_invoked}、\texttt{method}、\texttt{selected\_model} / \texttt{selected\_class}、\texttt{confidence}、\texttt{reasoning}、\texttt{candidates\_summary}（候选摘要）、\texttt{judge\_summary}（若启用 radiomics judge）、\texttt{llm\_model} / \texttt{temperature} / \texttt{max\_tokens}、时间戳 \texttt{ts}；可选 \texttt{prompts}（\texttt{system} / \texttt{user}）用于补充材料逐字复现。表 R11 中 \textbf{10} 例均来自 TN3K 主流程与 5× \textbf{分割 LLM 路由}稳定性实验的审计日志；修订稿将附 \textbf{≥10} 条去标识化样例记录（与表 R11 病例 ID 一致）。
\par
\begin{responsetable}
\caption{表 R12. Supplementary Table S[X] --- 基线训练与推理协议（完整）}
\label{tab:R12}
\begin{adjustbox}{max width=\textwidth}
\begin{tabular}{lcccc}
\toprule
模型 & 任务 & 甲状腺开发集 & 上游预训练 & 外测推理输入 \\
\midrule
BiomedCLIP & 分类 & dataset\_3 & OpenCLIP BiomedCLIP；在 dataset\_3\_cls 上微调分类头 & 仅图像 \\
LSNet & 分类 & dataset\_3 & ImageNet 预训练 LSNet & 仅图像 \\
MedSAM2 & 分割 & dataset\_4 & SAM 2.1 Hiera & 图像 + GT box/mask（\texttt{gt\_mask}） \\
MedSegX & 分割 & dataset\_4 & SAM ViT-B + MedSegX 甲状腺适配器 & 图像 + GT box/mask（\texttt{gt\_mask}） \\
MedSigLIP & 分类 & dataset\_3 & Google MedSigLIP；在 dataset\_3\_cls 上部分微调 & 仅图像 \\
RepViT & 分类 & dataset\_3 & ImageNet 预训练 RepViT & 仅图像 \\
ResNet-50 & 分类 & dataset\_3 & ImageNet ResNet-50；dataset\_3 微调 & 仅图像 \\
TransUNet & 分割 & dataset\_4 & ImageNet-21k ViT-B/16（TransUNet 默认） & 图像 + GT box/mask（\texttt{gt\_mask}） \\
UltraFedFM & 分类 & dataset\_3 & ViT 骨干 + dataset\_3\_cls 微调 & 仅图像 \\
UltraFedFM & 分割 & dataset\_4 & MAE / 联邦 FM + dataset\_4 微调 & 图像 + GT box/mask（\texttt{gt\_mask}） \\
ThyroidXAgent（分割路由） & 分割 & dataset\_4 & DINOv3-UNet（ImageNet 级骨干） & 图像 → LLM 路由后的专家 mask \\
ThyroidXAgent（分类路由） & 分类 & dataset\_3 & DINOv3 多任务骨干（dataset\_3 微调） & 路由 mask + 图像（pipeline） \\
GPT-5 / Gemini / MedGemma / Qwen3-VL（若见于 S4） & 分类 & ---（未微调） & 厂商/公开 VLM 权重 & 图像 ± 固定文本 prompt（零样本或 API） \\
\bottomrule
\end{tabular}
\end{adjustbox}
\responsetablenote{1. \textbf{dataset\_3（分类开发）}：TN3K + ThyroidXL + TN5K 堆叠训练划分（与 OpenThyroid 公开 multitask 分类开发集一致）。2. \textbf{dataset\_4（分割开发）}：TN3K + ThyroidXL + PKTN + TN5K 堆叠训练划分（与 OpenThyroid 公开分割开发集一致）。3. \textbf{监督基线默认协议}：在对应堆叠开发集上\textbf{全量微调}（自公开上游预训练权重出发），取验证集最佳权重用于 Table S3/S4 全部外测；\textbf{非}甲状腺零样本或仅用原始权重评估。4. \textbf{分割基线外测}：图像 + \textbf{GT box/mask} 提示（\texttt{gt\_mask} 协议），与 Table S3 一致。5. \textbf{分类基线外测}：仅超声图像输入。6. \textbf{ThyroidXAgent}：满堆叠 S3/S4 结果与上表使用相同堆叠开发数据；\textbf{Table S6} 逐级叠加消融另报。7. \textbf{视觉–语言模型}：未在 dataset\_3 端到端训练；厂商/公开权重；API 或冻结推理 + 固定 prompt。8. 监督行与 \textbf{Table S3/S4} 出表模型一致；ZJH-8K、DDTI 等外测队列\textbf{不参与}开发微调。}
\end{responsetable}

\begin{responsetable}
\caption{表 R13. 小于 5 mm 结节 --- 流水线失败统计}
\label{tab:R13}
\begin{adjustbox}{max width=\textwidth}
\begin{tabular}{lccccc}
\toprule
队列 & n 匹配 & <5.0 mm (N) & 失败 (M) & 小结节子集 mean Dice & GT 等价径 <50 px \\
\midrule
ddti\_seg & 637 & 54 & 1 & 0.734 & 54 \\
pktn & 150 & 15 & 1 & 0.741 & 15 \\
rjh7k & 7288 & 366 & 1 & 0.766 & 366 \\
thyroidxl & 2090 & 0 & 0 & --- & 0 \\
tn3k & 614 & 68 & 2 & 0.771 & 68 \\
tn5k & 1000 & 103 & 12 & 0.696 & 103 \\
zjh8k & 7958 & 1001 & 44 & 0.818 & 1001 \\
\bottomrule
\end{tabular}
\end{adjustbox}
\end{responsetable}

\subsection*{附录 A2. 放射组学 / AutoGluon 规范（回复信摘录）}
\label{app:A2}
\textbf{1. 管线位置}
\par
1. Phase 3 分割路由得到 \textbf{selected mask}。
\par
2. \textbf{无} radiomics 专用连通域后处理；特征提取与训练阶段使用同一套二维 PyRadiomics 设定。
\par
3. Phase 4：PyRadiomics 特征 → \textbf{AutoGluon-Tabular} 推断；恶性概率写入证据库，并供路由与分类综合使用。
\par
4. 可选：对表格模型计算 \textbf{SHAP} 特征归因（与稿面架构图一致）。
\par
\textbf{2. PyRadiomics}
\par
\begin{responsetable}
\caption{表 A2-1. PyRadiomics 特征提取设定}
\label{tab:A2-pyradiomics}
\begin{adjustbox}{max width=\textwidth}
\begin{tabular}{ll}
\toprule
项目 & 设定 \\
\midrule
软件版本 & \textbf{pyradiomics 3.1.0}（与训练环境一致） \\
参数设定 & 推理与训练共用同一套二维 PyRadiomics 配置（完整参数表见修订稿补充材料） \\
几何 & 强制 2D；mask 几何校正；重采样间距 \textbf{(1.0, 1.0)} 像素（各向同性假设） \\
强度 & 归一化 scale \textbf{100}；直方图 binWidth \textbf{25}；对称 GLCM \\
图像类型滤波 & Original、Wavelet、LoG（σ=2, 4）、Square、SquareRoot、Logarithm、Exponential、Gradient \\
特征类 & 完整 \textbf{shape2D}；\textbf{GLCM} 及预设纹理特征族 \\
输入 & 灰度超声图；mask 前景为阳性像素 \\
特征维度 & 训练提取 \textbf{100} 列 radiomics；AutoGluon 拟合时使用 \textbf{86} 个数值特征（内部丢弃 14 列稀疏 glszm） \\
\bottomrule
\end{tabular}
\end{adjustbox}
\end{responsetable}

\textbf{3. AutoGluon-Tabular}
\par
\begin{responsetable}
\caption{表 A2-2. AutoGluon-Tabular 训练与部署设定}
\label{tab:A2-autogluon}
\begin{adjustbox}{max width=\textwidth}
\begin{tabular}{ll}
\toprule
项目 & 设定 \\
\midrule
软件版本 & \textbf{1.4.0} \\
预设 & best\_quality \\
总时限 & \textbf{600 s}（DyStack + 主拟合） \\
Bagging & \textbf{8-fold}，\textbf{1} 层 stack \\
优化指标 & 二分类 \textbf{ROC-AUC} \\
训练样本 & \textbf{17,774} 张图像（良/恶性标签） \\
训练 mask & 堆叠 \textbf{dataset\_3} 上的 \textbf{GT mask}（与 Table S4 监督分类开发堆叠一致） \\
部署模型 & AutoGluon \textbf{加权集成}（bagged LightGBM、CatBoost、神经网络等；在开发集上自动选型与加权） \\
外测推理 & 固定上述开发集训练得到的集成模型，不在外测数据上更新权重 \\
\bottomrule
\end{tabular}
\end{adjustbox}
\end{responsetable}

\textbf{4. 与 R2.20 的区分（表注）}
\par
连通域阈值（专家内 ≥50 px、5×5 闭运算、集成 mask 取最大连通域）作用于 \textbf{DINO 专家/集成输出}，\textbf{不是} radiomics 前的额外 CCA 步骤。

\par
